\Article{32}
But what is the value of the angle $A$ when the attraction of the vessel is greater than the mutual attraction of the parts of the fluid? The imaginary value that for $\beta\beta>\alpha\alpha$ the formula $\sin\frac12 A=\frac\beta\alpha$ assigns to the angle $A$ already shows that some underlying assumption cannot be admitted in this case. Indeed, whenever $\beta\beta>\alpha\alpha$, the assumption that the surface $T$ has a \textit{boundary} cannot be reconciled with the condition that the function $W$ be a minimum. Wherever you place the boundary, imagine a very thin film of fluid spread out beyond it, so that $T$ increases by $T'$ and $U$ consequently increases by nearly the same amount. The value of the function $W$ clearly undergoes a change equal, up to an imperceptible difference, to the negative quantity $-(2\beta\beta-2\alpha\alpha)T'$. Indeed, the value of $W$ can still decrease further until $T'$ has occupied the whole remaining surface of the vessel. The value $-(2\beta\beta-2\alpha\alpha)T'$ for the change will be the more exact, the smaller we take the thickness to be. Insofar as only the value of the expression $W$ is at issue, nothing prevents us from imagining that the thickness decreases until it vanishes. A film of vanishing thickness, which must be carefully distinguished from one of imperceptible thickness, would, however, be a mere mathematical fiction. The shape of the region $s$ adapted to that fiction would in fact be the same as the one for which $W$ attains a minimum in the case $\beta\beta=\alpha\alpha$.

Matters are somewhat different in our physical problem, where such an additional film must have a definite thickness, however imperceptible, if equilibrium is to be possible. Whenever such a part is present, the expression $W$ is incomplete, as we showed in Article 18. If we denote by $T'$ the part of the vessel that the film covers,
\SourceLeaf{76}and by $\rho$ its thickness at any point, the terms to be added to the expression $\Omega$ are
\[
\pi cc\int\theta'\rho.\mathrm{d}T'-\pi cC\int\Theta'\rho.\mathrm{d}T'
\]
and hence the terms to be added to the value of $W$ are
\[
\begin{gathered}
\frac{\pi C}g\int\Theta'\rho.\mathrm{d}T'-\frac{\pi c}g\int\theta'\rho.\mathrm{d}T'\\
=\int\mathrm{d}T'.\left(\frac{2\beta\beta}{\Theta0}.\Theta'\rho-\frac{2\alpha\alpha}{\theta0}.\theta'\rho\right)
\end{gathered}
\]
Consequently, since the value of $W$ has already undergone the change $(2\beta\beta-2\alpha\alpha)T'$ through the addition of that film, the total change to be added to the value of $W$ obtained by leaving out the film will be
\[
-2\int\mathrm{d}T'.\left[\beta\beta\left(1-\frac{\Theta'\rho}{\Theta0}\right)-\alpha\alpha\left(1-\frac{\theta'\rho}{\theta0}\right)\right]
\]
Because $\theta'0=\theta0$ and $\Theta'0=\Theta0$, this change would be zero for vanishing thickness. But since $\theta'\rho$ and $\Theta'\rho$ decrease very rapidly as the thickness $\rho$ increases and are already imperceptible for an imperceptible value of $\rho$, the change will converge very rapidly toward the value $-(2\beta\beta-2\alpha\alpha)T'$. In the fluid’s state of equilibrium it will have to agree with that value up to an imperceptible difference, so that the value of the corrected expression $W$ cannot undergo a further perceptible decrease. A complete investigation of the law that the thickness $\rho$ must follow would require a deeper analysis. We shall not dwell on it here, since without knowledge of the functions $f,F$, on which the functions $\theta',\Theta'$ depend, and for the reasons to be given in Article 34, such an analysis might seem rather fruitless. To investigate the substantial part of the fluid, that is, the part whose dimensions are all perceptible, it is enough in our case, where $\beta\beta>\alpha\alpha$, to imagine the vessel \textit{wetted} near the boundary of the substantial part, that is, coated with a fluid film whose thickness is imperceptible yet large enough for $\theta'\rho$ and $\Theta'\rho$ to be neglected. The function that must be a minimum in the state of equilibrium will then be
\[
\int z\mathrm{d}s-2(\beta\beta-\alpha\alpha)(T+T')-\alpha\alpha T+\alpha\alpha U
\]
where $T,U$ are understood to refer only to the substantial part of the fluid. Clearly, then, the variation of this function resulting from a virtual change in the shape of the substantial part of the fluid, a change that leaves the sum $T+T'$ unaffected, agrees with the variation of the expression
\SourceLeaf{77}%
\[
\int z\mathrm{d}s-\alpha\alpha T+\alpha\alpha U
\]
that is, the very expression that must be a minimum in the case $\beta\beta=\alpha\alpha$. We conclude that the equilibrium shape of the fluid in a vessel for which $\beta\beta>\alpha\alpha$ agrees with the equilibrium shape of the same fluid in a vessel for which $\beta\beta=\alpha\alpha$, with the difference that in strict equilibrium the former must end in a film of imperceptible thickness. Moreover, the illustrious \textsc{Laplace} already observed that in the former case a vessel coated with a fluid film of imperceptible thickness is equivalent to a vessel whose particles exert an attractive force on the fluid equal to the mutual attractive force of the parts of the fluid.

This immediately gives the modification to be added to the propositions in Article 18 on the rise of fluids in vertical capillary tubes. Whenever $\beta\beta>\alpha\alpha$, we must substitute $\alpha\alpha$ for $\beta\beta$ in the formulas given there.

\Article{33}
In the case where $\beta\beta<\alpha\alpha$, the vessel cannot be wetted by a film of fluid of imperceptible thickness if the law of the functions $\theta',\Theta'$ is such that the value of the function
\[
\alpha\alpha\left(1-\frac{\theta'\rho}{\theta'0}\right)-\beta\beta\left(1-\frac{\Theta'\rho}{\Theta'0}\right)
\]
which for brevity we shall write as $Q\rho$, increases continuously as $\rho$ passes from the value $0$ toward a perceptible value. Clearly, with the function $Q\rho$ behaving in this way, the existence of such a film would conflict with the condition for a minimum. This behavior follows at once from the hypothesis we discussed in Article 31, in which $fx,Fx$ are in a fixed ratio independent of $x$. For it also follows from this that $\frac{\theta'\rho}{\theta'0}=\frac{\Theta'\rho}{\Theta'0}$ and hence $Q\rho=(\alpha\alpha-\beta\beta)(1-\frac{\theta'\rho}{\theta'0})$. But if the functions $f,F$ followed different laws, it would not be impossible for the value of $\frac{\theta'\rho}{\theta'0}$ to decrease more rapidly than that of $\frac{\Theta'\rho}{\Theta'0}$ and for the function $Q\rho$, within the range of imperceptible values of $\rho$, first to become negative and then, after reaching its minimum value (that is, its negative extreme), to rise again through the value $0$ toward its positive limit $\alpha\alpha-\beta\beta$. In such a case equilibrium would indeed require an imperceptible film whose thickness, generally speaking, would have to be such that $Q\rho$ does not differ perceptibly from its minimum value. If we denote that minimum by $-\beta'\beta'$, we shall have $\beta'\beta'<\beta\beta$. The shape of the substantial part of the fluid will be determined just as if it were in a vessel for which we must replace the quantity $\beta\beta$ with $\beta'\beta'$. That is, the angle between the plane
\SourceLeaf{78}tangent to the free surface of the fluid at the boundary of the substantial part and the wall of the vessel will be $=2\operatorname{arc.}\sin\frac{\beta'}{\alpha'}$. But since it is very doubtful whether such a case exists in nature, it seems needless to dwell on it any longer.

